\section{TF chain validation as a graph search}
\label{sec:tf}

ROS transforms form a tree: each frame has exactly one parent. Localization is
healthy only if the chains the stack needs, \code{map}$\to$\code{odom} and
\code{odom}$\to$\code{base\_link}, exist and are fresh.

Treat every recently received transform as a directed edge
parent$\to$child; static transforms never go stale. For each required chain
$(s, t)$, a depth-first search from $s$ looks for $t$. Three things can be wrong:
\begin{enumerate}[nosep]
  \item \emph{No path}: a node stopped publishing (MISSING).
  \item \emph{Two parents}: two nodes publish the same child, for example both
        the EKF and the simulator publishing \code{odom}$\to$\code{base\_link}.
        Edges are keyed by publisher identity, so this shows up as a frame with
        two parents on the path (WARN).
  \item \emph{A cycle}: a search that revisits a frame on its own path (WARN).
\end{enumerate}
Each search visits every edge at most once, so the check is linear in the
number of transforms.
\file{mhseals\_testing/monitor.py (\_tf\_rows)}
